%% !TEX root = ../manual_gaolv5.tex
% The introduction (unnumbered).
%%
%% Copyright (c) 2001 Swiss Federal Institute of Technology, Switzerland
%% Copyright (c) 2002-2016 Laboratoire d'Informatique de
%%                         Nantes-Atlantique, France
%% Copyright (c) 2017-2020 Laboratoire des Sciences du Numérique de Nantes,
%%                         France
%% Copyright (c) 2026 ENSTA, France
%%
%% By: Frederic Goualard, Jordan NININ
%%
%% GAOL is a software distributed WITHOUT ANY WARRANTY. Read the associated
%% COPYING.LIB file for information.
\chapter*{Introduction}\label{chap:introduction}
\addcontentsline{toc}{chapter}{Introduction}
\thispagestyle{plain}

Computers compute with floating-point numbers, and a floating-point number
is rarely the real number a program means: $0.1$ has no finite binary
expansion, the quotient $1/3$ is rounded, and the rounding errors of a long
computation add up in ways that are hard to follow. Interval arithmetic
answers this difficulty by computing with pairs of floating-point numbers:
an \emph{interval} $[a, b]$ stands for every real number between $a$ and
$b$, and each operation on intervals returns an interval that contains all
the results the operation can take on the reals it encloses. A program that
computes with intervals proves, at the end of its computation, that the
exact result lies within the interval it obtained. This is the containment
property, which the first chapter defines, and which every operation of GAOL
honours.

\medskip
GAOL\footnote{For those readers who are not native English speakers,
``gaol'' should be pronounced \textsf{j\={a}l}, like the word ``jail'', of
which it is a, chiefly British, variant.}
\cindex{gaol!pronunciation}%
is a C++ library to perform arithmetic with floating-point intervals. Its
development was initiated at the \emph{Swiss Federal Institute of
Technology}, Lausanne, Switzerland, while Fr\'ed\'eric Goualard was a
postdoctoral fellow supported by the \emph{Swiss National Science
Foundation}. It started as a limited version of \emph{Jail} (now
\emph{halloween}), a templated C++ interval library developed during
Goualard's PhD. To our knowledge, a unique feature of GAOL among the C++
interval arithmetic libraries available when it was written is the
implementation of \emph{relational arithmetic operators}, which interval
constraint programming requires (Chapter~\ref{chap:arithmetic}).
\cindex{operator!relational}\cindex{relational operator}%
Hence the pun of the name: GAOL is \emph{not} JAIL (\emph{Just Another
Interval Library}). The standard for interval arithmetic, IEEE
1788-2015~\cite{IEEE1788:2015}, has since required such operators, as its
\emph{reverse functions}.
\cindex{IEEE 1788-2015}%

\medskip
GAOL v5, written by Jordan Ninin at ENSTA, continues GAOL as Fr\'ed\'eric
Goualard developed it, from its version 4.2.2. It was written for Codac and
IBEX, whose intervals are built upon GAOL, and brings the changes GAOL needs to
compile and compute right on every system it can: Linux, macOS and Windows,
on x86, x86-64, 32-bit and 64-bit ARM, POWER (ppc64le), s390x and RISC-V
(riscv64).
\cindex{GAOL v5}%
What it brings, and this manual describes, is the following.

\begin{itemize}
\item[$\bullet$] \textbf{Correctly rounded elementary functions.} Every
  elementary function that CORE-MATH~\cite{CoreMath:2022} provides is bounded
  with it, which is correctly rounded in the rounding direction in effect: the
  value GAOL computes at a bound, rounding upward, \emph{is} the upper bound,
  and the double below it the lower one, the tightest bounds. The integer
  powers, the roots of an order other than 2 and 3 and \code{compoundm1}, which GAOL computes itself, are the
  tightest bounds or within a few doubles of them
  (Chapter~\ref{chap:accuracy}). The bounds are the same on every architecture
  and with every compiler, and there is no mathematical library to install, to
  find or to link along with GAOL: CORE-MATH is compiled into it.
\item[$\bullet$] \textbf{IEEE 1788-2015.} GAOL v5 follows the standard for
  interval arithmetic in the operations GAOL provides: the empty set, the
  infinite bounds, the domains of the functions, the comparisons and the
  interval literals. The few cases where it differs, mostly the \code{pow}
  of the namespace \code{gaol}, are stated where the operations are
  described, and the tightness of each operation is documented, as the
  standard requires (Chapter~\ref{chap:accuracy}).
\item[$\bullet$] \textbf{New operations}, which the standard requires or
  recommends: \code{exp2}, \code{exp10}, \code{log2}, \code{log10},
  \code{fma}, the integer functions \code{sign}, \code{trunc},
  \code{round\_ties\_to\_even} and \code{round\_ties\_to\_away}, sixteen
  recommended functions (\code{expm1}, \code{log1p}, \code{hypot},
  \code{compoundm1}, \code{sinpi}\ldots), the cancellative subtraction and
  addition, the radius, the comparisons \code{less} and
  \code{strictly\_less}, the boolean functions \code{is\_entire} and
  \code{is\_common\_interval}, \code{overlap()}, the reverse functions of the
  power and of \code{atan2} (\ieee{powRev1}, \ieee{powRev2}, \ieee{atan2Rev1},
  \ieee{atan2Rev2}), \ieee{pownRev} for a non-positive exponent and
  \ieee{mulRevToPair}; \code{atan2}, which GAOL declared, is implemented.
\item[$\bullet$] \textbf{The names of the standard.} The namespace
  \code{gaol\_ieee1788} gives each operation of the standard that GAOL
  provides the name and the argument order the standard gives it
  (Chapter~\ref{chap:ieee1788}).
\item[$\bullet$] \textbf{The rounding direction is the business of GAOL.}
  Each operation that computes a bound sets the rounding direction upward
  when it is not, whichever direction the calling code left; the
  operations of the AVX-512 path set none, their instructions carrying the
  direction. The user no longer has to keep it upward
  (Chapter~\ref{chap:rounding}).
\item[$\bullet$] \textbf{Exact input and output.} Numbers are read exactly,
  and intervals written in decimal are rounded outward, whatever the C
  library; written in hexadecimal, an interval is read back as the same
  doubles (Chapter~\ref{chap:io}).
\item[$\bullet$] \textbf{Three builds that agree.} GAOL has a CMake build,
  which is the one to use, and the autotools and meson builds, and the three
  configure GAOL the same way. Tests compare the bounds GAOL computes with
  their exact values, and a continuous integration builds and tests GAOL on
  Linux, macOS and Windows, with the compilers and on the processors listed in
  Chapter~\ref{chap:installation}.
\end{itemize}

\noindent
A program written for GAOL~4 needs a few changes, which
Chapter~\ref{chap:changes} lists: first of all, \file{gaol/gaol.h} no longer
opens the namespace \code{gaol}.

\section*{About this manual}

This document is both a manual and a reference to use GAOL. It assumes a
prior knowledge of interval arithmetic, whose notions the first chapter
recalls in the form GAOL uses them; the books and papers by Goldberg,
Neumaier, Moore, Alefeld and Herzberger, and
others~\cite{Goldberg:01,IEEE754:2008,Moore:Book1966,Alefeld-Herzberger:83,Neumaier:Book1990}
give a full presentation of floating-point arithmetic, interval arithmetic
and the use thereof.
\cindex{references!on interval arithmetic}%
The main entry point for interval arithmetic on the Web is the \emph{Interval
Computation site} of Vladik Kreino\-vich,
\url{http://www.cs.utep.edu/interval-comp/}.
\cindex{interval!computation site (web)}

\medskip
The pages of the directory \file{doc} of the sources give the details of each
change GAOL v5 brings, with its measures, and this manual takes their content
up; Appendix~\ref{chap:documentation} lists them. The mark
\textcolor{blue}{\textsf{\scriptsize[GAOL~v5]}} follows, throughout the
text, what behaves otherwise than in GAOL~4. A clause number in parentheses,
as (12.10.3), and a table cited as Table~9.1 or Table~10.7 are those of
IEEE 1788-2015 (the tables of this manual are numbered by chapter and linked).

\medskip
The examples of the reference show what they print, and every output shown
was produced by GAOL v5. A script of the sources compiles each example that
shows its output under its code, runs it and compares what it prints with the
manual (Chapter~\ref{chap:tests}); the outputs shown apart, in a frame
Output, are not compared by it. An
example that writes a \code{bool} sets \code{boolalpha} first: without it,
\code{cout} writes 1 and 0. The examples assume \code{using namespace gaol;}
and \code{using namespace std;}, except those of
Chapter~\ref{chap:ieee1788}, which open \code{gaol\_ieee1788} instead, and
each example is a program of its own.

\medskip
Classes, methods, functions, macros, constants and variables available in
the library but not described in this document are likely to change or to be
removed. Consequently, they should be used with caution, if at all. Among
them are the intervals of floats, \code{gaol::intervalf} and
\code{gaol::interval2f}, which are unfinished and which none of the builds
compiles: only a developer of GAOL does, defining
\code{GAOL\_FLOAT\_INTERVALS} in \file{gaol/gaol\_config.h}.
\cindex{undocumented functions}%
\cindex{function!undocumented}%

\section*{Outline}

The first chapter recalls the notions this manual relies on: floating-point
numbers and their roundings, intervals and the containment property, the
dependency problem, what IEEE 1788-2015 fixes, the idea of the relational
functions, the mathematical libraries from which correctly rounded
elementary functions came, and the accuracy modes of the standard.

\medskip
Chapters~\ref{chap:installation} and~\ref{chap:overview} are the user's
guide. The second chapter describes the installation of GAOL: the platforms
it supports, the compilers and options it refuses, the three builds and their
options, and what the configuration decides by itself. The third chapter
shows how a program uses GAOL: a first program, the flags of interval
arithmetic the code including the headers of GAOL is compiled with, the two
namespaces and their powers, the tools of interval algorithms, the common
errors, and the correspondence with the names of other interval libraries.

\medskip
Chapter~\ref{chap:rounding} is devoted to the rounding direction and to the
floating-point environment, which GAOL v5 takes charge of: what the upward
rounding does to the rest of the program, how the direction can be
preserved, the AVX-512 path, the NEON intervals of 64-bit ARM processors, the
FPU intervals, the modes that flush the subnormal numbers to zero, the
floating-point exceptions, and the initialization and the cleanup of GAOL.

\medskip
Chapters~\ref{chap:creation} to~\ref{chap:io} are the reference of the class
\code{interval}: the constructors and the assignments, the relations, the
arithmetic operations, functional and relational, the numeric and integer
functions, the splitting methods and the set operations, and the reading and
writing of intervals. Chapter~\ref{chap:ieee1788} gives the names of IEEE
1788-2015, and where the standard and GAOL differ. Chapter~\ref{chap:accuracy}
states the tightness of each operation, as the standard requires.

\medskip
Chapter~\ref{chap:computing} explains how GAOL computes its intervals: the
representation of an interval in a register, the SSE2, AVX-512, NEON and FPU
kernels, where the time of a program goes, the cost of the rounding direction
on each processor, what the compilers give, how GAOL compares with the other
interval libraries, with ITF1788, the test suite of the standard, and with the
benchmark of Tang et al., and the verdicts
on new paths. Chapter~\ref{chap:errors} describes the floating-point constants
and functions, the version information, the errors and the debugging
facilities. Chapter~\ref{chap:tests} presents the tests, the examples, the
continuous integration and the scripts of the repository.

\medskip
The last chapter, Chapter~\ref{chap:changes}, is written for the programs
written for GAOL~4: what they have to change to compile with GAOL v5, and what they will
see differently. The conclusion sums up what GAOL v5 established and what is
left open. The bibliography comes before two appendices, which list the
documentation and the references of GAOL, with the way to report a bug, and
reproduce the licences of GAOL and of CORE-MATH; an index ends the manual.
